\section{模型的检验、灵敏度分析与评价}\label{sec:eval}

\subsection{正确性检验的内容}

本文的每一次实验都包含逐例核验：清除数是否等于真实源数；每次成功清除时清除点到真源的距离是否不超过 $20\mmet$；每一时刻的可行多边形与其最小包围圆是否确实包含真源（即\cref{eq:invariant}）；逐条动作重新累加的时间与日志时钟之差是否在 $10^{-6}\msec$ 以内；测站布局是否在连续域上满足覆盖条件；各类边界与退化几何是否处理正确；候选版本有无退步案例；以及失败运行是否被独立保留。

其中第三项尤其关键。\cref{app:debug}记录的那次失败正是“清除数看似完全正确、几何证书却已经失效”的情形，它说明\textbf{结果正确不意味着不做检查}。

\subsection{灵敏度分析}

\paragraph{对测向误差界的灵敏度.}

改变实际误差幅度、同时把算法的误差界同步放宽为“幅度 $+0.005^\circ$”，每档使用相同的种子，结果见附录~\ref{app:sens}的\cref{tab:senserr}。该组实验做在早期的“两次测向”基线上，只用于观察趋势，其绝对数值不代表当前版本。

误差幅度加倍只带来 $2\%$（问题三）与 $5\%$（问题四）的时间增长，说明在误差界被正确告知的前提下，方案对误差幅度不敏感。原因是：可行域的尺度与 $\bar\varepsilon$ 成正比，但总时间的主体是与 $\bar\varepsilon$ 无关的巡站与扫描固定成本。

\paragraph{对误差的失配是危险的.} 需要强调的是，这种“不敏感”并不意味着方案对误差界的设定可以放松；相反，误差界的失配可能导致可行域被错误裁剪。附录~\ref{app:sens}的表~\ref{tab:senserr}只说明“正确知道更大误差界”时的灵敏度。本文专门构造了真实偏差达 $1.5^\circ$ 而算法只容许 $1.005^\circ$ 的算例，结果出现正的半平面违约 $12.958908\mmet$——真源被排出了可行域，包含性不变量被破坏。这正是算法采用略宽于 $1^\circ$ 的保守界 $\bar\varepsilon=1.005^\circ$ 的原因：吸收“误差先产生、再按两位小数取整”可能带来的额外偏差。一旦正式测试中出现空交集或“认证清除失败”，正确的做法是排查取整、角度约定、坐标与计时，确认无误后再决定是否调整阈值。因此，在误差界能够被正确估计并同步调整的前提下，方案对误差幅度变化总体不敏感。

\paragraph{对最小接收半径的灵敏度.} 这是全文最关键的敏感参数，因为两问的覆盖保证都以 $\rho_{\min}=1000\mmet$ 为前提。第三问相当宽裕：由\cref{thm:hexcover}，七站布局可行当且仅当 $\rho_{\min}\ge R/2=900\mmet$，题设 $1000\mmet$ 有 $11.1\%$ 余量。第四问则薄得多：$22$ 站布局在 $\rho_{\min}\ge994\mmet$ 时仍通过、$990\mmet$ 时失效（附录~\ref{app:sens}表~\ref{tab:sensrho}），容忍度约 $6\sim10\mmet$；表中 $998\sim1000\mmet$ 的余量不变、$994\mmet$ 突降是粗网格量化所致，题设下的真实几何余量以 $h=1\mmet$ 细网格的 $1.7331\mmet$ 为准（\cref{sec:q4:cert}）。值得注意的是，优化搜索得到的 $22$ 站布局不仅比解析的 $25$ 站更快，对 $\rho_{\min}$ 的余量也更大（约 $6\sim10\mmet$ 对不足 $2\mmet$）；两者余量都不宽，故程序保留回退开关：删除配置中的 \texttt{q4\_station\_list} 即恢复 $25$ 站布局。

\paragraph{对光学格距与其他阈值的灵敏度.} 光学兜底要求 $h_{\rm opt}\le28.2843\mmet$（取 $25\mmet$），条带覆盖要求外接圆半径 $\le r_c$（取 $19.9\mmet$），认证阈值取 $20\mmet-10^{-6}$。这些阈值一旦被误改超限，覆盖保证立即失效，因此程序启动时统一校验、不满足即拒绝运行。

\subsection{稳健性（压力场景）}

除随机案例外，本文设计了五类压力场景：源全部贴近区域边界且取最小接收半径、源在圆心聚集、全体取最小接收半径、第四问全部为朝外的定向源、固定 $\pm1^\circ$ 偏差与光滑空间相关偏差。\cref{tab:q4holdout}给出了第四问五类场景的完整结果：$220$ 例全部清除，$22$ 站在每一类上都优于 $25$ 站。第三问的压力集（$70$ 例、$890$ 源）同样全部清除，只有 $2$ 例轻微退步，最大退步 $2.72\msec/$源。

\subsection{模型的优点}

\textbf{核心贡献有四点。第一，“一定能清完”是一条可证明的性质}：\cref{thm:term3}与\cref{thm:hull}把终止条件建立在连续域几何覆盖上，\cref{thm:cert}把它化为可执行的有限检查，\cref{cor:certify}与\cref{eq:grid}、\cref{eq:strip}保证清除动作确定性成功，整条链路上没有一步靠“试几次收不到就认为没有”。\textbf{第二，七站达到距离覆盖的理论下界}：\cref{thm:seven}证明 $6$ 站不可能、$7$ 站可达。\textbf{第三，凸包条件刻画定向源的方向覆盖，$22$ 站取得连续域证书}：该布局在 $1018.68$ 万点细网格上最小余量 $1.7331\mmet$，所比较布局中 $\le21$ 站均无法通过。\textbf{第四，时间诊断驱动布局优化}：先由\cref{sec:analysis:fixed}认定固定成本主导，再把优化落到测站布局上。此外，题目只给出误差的界，本文据此只用凸多边形交与最小包围圆，并对所有近似取保守外包，使包含性不变量单调成立。

\textbf{工程优化另有若干}：\texttt{no\_signal} 圆盘排除、认证条带覆盖、共享测站扫描、圆心补测精化、有界等待与即时清除等，贯穿第\ref{sec:q3}、\ref{sec:q4}节。附录~\ref{app:ablation}的表~\ref{tab:q3reject}与表~\ref{tab:q4reject}保留了 $11$ 条被否决的路线与一次几何证书失效，它们既是消融实验，也是结论可信度的一部分；覆盖保证参数在程序启动时统一校验，不满足即拒绝运行。

\subsection{模型的缺点与局限}

\begin{enumerate}[label=(\arabic*),itemsep=1pt]
  \item \textbf{调度层是启发式的.} 混合贪心只比较一步的移动代价，没有纳入未来的测量结果与光学路径；两步前瞻的尝试反而更慢（\cref{tab:q3reject}），但这不说明贪心已最优。
  \item 第四问的 \texttt{no\_signal} 尚未被利用. 一次“收不到”蕴含“距离超限或背向”的析取信息，可在“位置 $\times$ 朝向”的联合状态空间上做保守排除；本文原型未能胜过现行方案，这是最有价值的未完成方向。
  \item \textbf{覆盖余量偏薄.} $22$ 站布局对 $\rho_{\min}$ 只有约 $6\sim10\mmet$ 的余量（\cref{tab:sensrho}）；若官方对 $\rho_{\min}$ 的保证比 $1000\mmet$ 更弱，覆盖可能失效。
  \item \textbf{本地分布未必与官方一致.} \cref{tab:q4match}显示偏差随源数从 $-0.4\%$ 升到 $+8.4\%$，故本地配对实验的相对结论可谨慎外推、绝对数值不能。
  \item \textbf{异常只能安全停止.} 客户端会检测超时与协议异常、记录现场并停在原地退出，但没有自动恢复能力。
  \item \textbf{结论依赖题设的理想化物理.} 静止源、无遮挡直线行走、确定性阈值、误差有界，任一被打破，包含性不变量与覆盖证明都需重建。
\end{enumerate}

\subsection{模型的改进与推广}

针对上述局限，最值得推进的是把 \texttt{no\_signal} 的析取信息纳入“位置 $\times$ 朝向”的联合证书，并在多机协同下进一步摊薄固定成本。本文框架也可推广：区域改变时只需把\cref{thm:cert}的网格证书套用到新区域重新搜索布局；定向覆盖角收窄为 $2\beta$ 时把凸包条件换成更强的“任意 $\beta$-锥内至少有一站”；多机协同则化为测站集划分的 min--max $k$-TSP。水下声源搜寻、林区火点定位等“方位有界误差、作用距离有下界、处置有作用半径”的任务也可整体移植。

最后需要说明：本文的工程价值不仅在于给出一个可运行的方案，更在于区分了严格保证、数值验证与经验估计三类结论。
